\chapter{O REPL}

\section{O que é o REPL}

REPL (\textit{Read-Eval-Print Loop}) é o modo interativo de \hayashi{}. Cada linha
digitada é avaliada imediatamente e o resultado é exibido. É o ambiente ideal para
explorar dados, testar expressões e prototipar análises.

\begin{lstlisting}[style=inline, language=sh]
hay
\end{lstlisting}

\begin{lstlisting}[style=inline]
Hayashi 0.2.8-dev  — Applied Econometrics Language
In honor of Fumio Hayashi. Type 'exit' or Ctrl-D to quit.

hay>
\end{lstlisting}

\section{Saindo do REPL}

\begin{center}
\begin{tabular}{ll}
\toprule
\textbf{Comando} & \textbf{Efeito} \\
\midrule
\texttt{exit}   & Sai do REPL \\
\texttt{Ctrl-D} & Sai do REPL (EOF) \\
\texttt{Ctrl-C}                & Cancela a linha atual \\
\bottomrule
\end{tabular}
\end{center}

\section{Navegação e histórico}

O REPL usa \texttt{rustyline} e suporta:

\begin{itemize}
  \item Seta para cima/baixo: histórico de comandos
  \item \texttt{Ctrl-R}: busca no histórico
  \item \texttt{Ctrl-A} / \texttt{Ctrl-E}: início/fim da linha
  \item \texttt{Ctrl-C}: cancela a linha atual
  \item \texttt{Ctrl-D}: sai do REPL (equivalente a \texttt{exit})
\end{itemize}

O histórico é persistido entre sessões em \texttt{\$HOME/.hay\_history}.

\section{Expressões no REPL}

Qualquer expressão válida pode ser digitada diretamente:

\begin{lstlisting}[caption={Sessão interativa básica}]
hay> 2 + 2
4

hay> let x = 10
hay> x * 3
30

hay> let s = "olá, mundo"
hay> s
"olá, mundo"
\end{lstlisting}

\section{Executando scripts}

Scripts \texttt{.hay} são arquivos de texto com uma instrução por linha. Para executar:

\begin{lstlisting}[style=inline, language=sh]
hay analise.hay
\end{lstlisting}

A saída vai para \texttt{stdout}. Erros vão para \texttt{stderr}.

\section{Formato de erros}

\hayashi{} exibe erros com contexto visual:

\begin{lstlisting}[style=inline]
error: line 3: undefined variable 'y' — did you mean 'x'?
   3 │ display y + 1
     │ ^^^^^^^^^^^^^
\end{lstlisting}

O trecho de código, a linha e uma sugestão (quando disponível) facilitam a correção
sem precisar contar linhas manualmente.
